\section{Perceived AI Support}
\label{sec:eval-ai-support}

In addition to the SUS and R-TLX, after each mode participants answered a
six-item Perceived AI Support instrument on a five-point Likert scale. Item
$5$ (legal-basis confidence) is the primary confidence measure. The remaining
five items capture distinct facets of perceived AI support and are reported
descriptively here; the six-item composite is reported as a secondary
attitudinal measure.

\subsection{Per-Item Results}

Table~\ref{tab:ai-per-item} reports the mean response per item per mode and
the aggregated mean. Figure~\ref{fig:ai-subquestions} shows the corresponding
distributions.

\begin{table}[htbp]
\centering
\caption{Perceived AI Support, mean response per item by mode ($n = 69$ per
mode), five-point Likert scale (higher is better). The aggregated ``All''
column is the unweighted mean of the three per-mode means.}
\label{tab:ai-per-item}
\small
\begin{tabularx}{\linewidth}{@{}c X r r r r@{}}
\toprule
\# & Item & Form & Ask & Chat & All \\
\midrule
1 & The AI was helpful in creating this document. & 4.23 & 4.32 & 4.13 & 4.23 \\
2 & The AI helped me to complete the document faster. & 4.43 & 4.25 & 3.99 & 4.22 \\
3 & I understood why the AI suggested specific GDPR fields. & 3.51 & 4.16 & 3.75 & 3.81 \\
4 & I trust that the AI made correct suggestions. & 3.41 & 3.80 & 3.68 & 3.63 \\
5 & I am confident that the AI identified the correct legal basis. & 3.58 & 3.74 & 3.81 & 3.71 \\
6 & I would use this mode again. & 3.99 & 3.55 & 3.42 & 3.65 \\
\bottomrule
\end{tabularx}
\end{table}

\begin{figure}[htbp]
\centering
\includegraphics[width=\linewidth]{ai_subquestions_boxplot.png}
\caption{Box plots of the six Perceived AI Support items by mode
($n = 69$ per mode).}
\label{fig:ai-subquestions}
\end{figure}

For \emph{``The AI was helpful in creating this document''} (item~1),
Form scores $4.23$, Ask $4.32$, Chat $4.13$; aggregated $4.23$. Helpfulness
is rated highest in Ask and lowest in Chat, with a small absolute spread of
roughly $0.20$~points.

For \emph{``The AI helped me to complete the document faster''} (item~2),
Form is rated highest ($4.43$), Ask second ($4.25$), Chat lowest ($3.99$);
aggregated $4.22$. Perceived speed-up decreases monotonically with increasing
LLM involvement, dropping by approximately $0.45$~points from Form to Chat.

For \emph{``I understood why the AI suggested specific GDPR fields''}
(item~3), Form scores lowest ($3.51$), Chat second ($3.75$), Ask highest
($4.16$); aggregated $3.81$. Transparency is rated approximately
$0.65$~points higher in Ask than in Form. Ask is the only mode that
explicitly explains each suggestion.

For \emph{``I trust that the AI made correct suggestions''} (item~4), Form
scores lowest ($3.41$), Chat second ($3.68$), Ask highest ($3.80$); aggregated
$3.63$. Trust increases by approximately $0.27$--$0.39$~points from Form to
the LLM-driven modes.

For \emph{``I am confident that the AI identified the correct legal basis''}
(item~5, primary confidence measure), the three modes show a monotonic
increase: Form $3.58$, Ask $3.74$, Chat $3.81$; aggregated $3.71$. Confidence
in legal correctness rises by approximately $0.23$~points from Form to Chat.

For \emph{``I would use this mode again''} (item~6), Form scores highest
($3.99$), Ask second ($3.55$), Chat lowest ($3.42$); aggregated $3.65$.
Re-use intent decreases by approximately $0.57$~points from Form to Chat,
the largest negative gap of the six items.

\subsection{Composite Perceived AI Support}

The unweighted mean of all six items yields a composite Perceived AI Support
score of $3.86$ for Form, $3.97$ for Ask, and $3.80$ for Chat (aggregated
$3.87$). The omnibus Friedman test on the composite is non-significant
($\chi^2(2) = 1.05$, $p = 0.592$, $W = 0.008$), indicating that overall
perceived support does not differ meaningfully across modes once the six
items are pooled. The single legal-basis confidence item ($5_{ai}$) is also
non-significant under Friedman ($\chi^2(2) = 3.76$, $p = 0.152$). The re-use
item ($6_{ai}$) does reach significance ($\chi^2(2) = 11.12$, $p = 0.004$,
$W = 0.081$), with a large rank-biserial effect for Form versus Chat
($r = +0.543$). Full statistics are reproduced in
Appendix~\ref{app:inferential}.

\subsection{Mode Preference Ranking}

After all three modes had been completed, participants ranked the three modes
in order of overall preference ($1 =$ most preferred, $3 =$ least preferred).
Table~\ref{tab:mode-preference} summarises the ranking.

\begin{table}[htbp]
\centering
\caption{Post-experiment mode preference ranking ($n = 69$). ``Rank-1 votes''
counts the number of participants who placed the mode first.}
\label{tab:mode-preference}
\begin{tabular}{@{}l r r r@{}}
\toprule
Mode & Rank-1 votes & Share & Mean rank \\
\midrule
Form & 37 & 53.6\,\% & 1.78 \\
Chat & 20 & 29.0\,\% & 2.07 \\
Ask  & 12 & 17.4\,\% & 2.14 \\
\bottomrule
\end{tabular}
\end{table}

Form receives the plurality of first-place votes ($53.6\,\%$), followed by
Chat ($29.0\,\%$) and Ask ($17.4\,\%$). The mean ranks follow the same
ordering. A chi-square goodness-of-fit test against a uniform null rejects
the null hypothesis ($\chi^2(2, N = 69) = 14.17$, $p < 0.001$, Cramér's
$V = 0.320$, medium effect). The complementary Friedman test on the full
rank distribution is non-significant ($\chi^2(2) = 5.07$, $p = 0.079$),
reflecting that the strong Form plurality at rank~1 is partially balanced by
variance at ranks~2 and~3.

\subsection{Free-Text Comments}

Each mode collected a per-mode open-ended comment, and after all three modes
participants wrote an overall summary comment. A total of $88$ comments were
recorded across $46$ unique participants. Comments are written in German with
a small minority in English; below they are summarised thematically. The
section preserves the order Form, Ask, Chat, and concludes with the
across-modes overall comments.

\paragraph{Form.} Comments on Form cluster around speed, clarity, and
control. Recurring observations describe the mode as the fastest of the
three, the easiest to obtain an overview of, and the most ``integrated''
into the document creation workflow because the checkbox interface makes the
selection space visible at a glance. Several participants emphasise that
Form ``takes the most work off you'' (\emph{``nimmt einem die meiste Arbeit
ab''}) by pre-populating the option set, while a smaller subset notes the
opposite concern: that the absence of an interactive AI commentary makes it
unclear on what basis the model produces its suggestions, and that
participants would not feel safe relying on Form output without an additional
verification step. Several comments describe specific UX issues: long text
fields rendered as a single narrow line on mobile, the inability to provide
the AI with additional context for its suggestions, and the impression that
the AI ``guesses'' without scenario-specific input.

\paragraph{Ask.} Comments on Ask are the most polarised of the three modes.
A consistent theme is that Ask is the most cumbersome mode in terms of
workflow: the requirement to manually copy the AI's suggested text into the
form field is repeatedly described as redundant work, particularly because
participants reported that the model often produced suggestions in a
free-text style that did not match the structured form fields. Several
participants observed that the lack of a persistent context across fields
forced them to re-enter the same information repeatedly. Conversely,
explanatory text generated by Ask is consistently described as the most
informative of the three modes, supporting the high transparency rating
visible in item~3 of the AI Support instrument. A frequent suggestion is
to combine Ask's explanations with Form's structured field layout.

\paragraph{Chat.} Comments on Chat are dominated by two themes: comfort and
opacity. On the comfort side, several participants describe Chat as the most
intuitive of the three modes, especially because the AI takes over the work
of mapping the conversation onto the structured ROPA fields, leaving the
participant only with the task of describing the processing scenario. On
the opacity side, an equally large group of participants reports
uncertainty about what is actually being written into the document — whether
the entire chat history is used for generation, whether each individual
message must be confirmed, and what the document looks like during the
conversation. Specific defects mentioned include the AI occasionally
switching from German to English mid-response, occasional emission of raw
\texttt{DATA\_UPDATE} tokens, and a slow ``streaming'' rendering style that
participants experienced as artificially throttled.

\paragraph{Across modes (overall comments).} The overall comments
recapitulate the per-mode pattern. A clear majority describes Form as the
fastest and most ergonomic mode, particularly for participants who feel they
``do not know enough'' to interrogate the model meaningfully. Chat is
typically placed second and motivated by the impression that ``the AI takes
care of it,'' which several participants explicitly note increases their
\emph{subjective} sense of safety even though they cannot verify the
output. Ask is most often placed third, with the manual copy-paste step
cited as the principal reason. A recurring constructive suggestion across
participants is a hybrid mode that combines the visible option set of Form
with the conversational, context-aware behaviour of Chat or Ask. Several
participants explicitly add a caveat that they are novices in the legal
domain and that their preference ordering might shift if they had prior
GDPR knowledge.
